\chapter{平面曲线}\label{ch:b1:curves}

 \cref{ch:b1:derivative,ch:b1:taylor} 中的微\omterm{thm:b1:integration:def}{积分}，原是为
函数 $y = f(x)$ 而建立的；几何与力学中的大多数曲线——
例如轨迹、圆在圆上滚动所成的曲线、轨道——都无法写成这种形式，
而是以\emph{\omterm{def:b1:curves:def}{参数曲线}}的形式出现，形如 $t \mapsto (x(t),
y(t))$ 或以\emph{极\omterm{prop:b1:vspaces:coordinates}{坐标}}曲线 $r = r(\theta)$ 的形式出现。
本章研究这两类曲线，并以\omterm{def:b1:curves:conics}{圆锥曲线}作结。

\section{参数曲线}

\begin{definition}\label{def:b1:curves:def}
\emph{参数曲线}\index{参数曲线}是\omterm{def:b1:logic:map}{映射} $f
\colon I \to \R^2$，$t \mapsto M(t) = (x(t), y(t))$，其中 $x, y$ 为
（至少）$C^1$ 类、定义于\omterm{prop:b1:reals:intervals}{区间} $I$ 上的函数。\emph{速度向量}
为 $f'(t) = (x'(t), y'(t))$；点 $M(t_0)$ 在
$f'(t_0) \neq (0,0)$ 时称为\emph{正则点}，其在该点处的\emph{切线}
为过 $M(t_0)$ 且以 $f'(t_0)$ 为方向的直线。
\end{definition}

\begin{proof}[为何 $f'$ 决定切线方向]
逐分量应用泰勒–杨公式，$M(t) = M(t_0) + (t - t_0) f'(t_0) +
o(t - t_0)$：弦的方向 $\frac{M(t) - M(t_0)}{t - t_0}$
趋于 $f'(t_0)$。正则性假设正是使
该极限成为\emph{方向}：若 $f'(t_0) = 0$，则上式退化为
$M(t) = M(t_0) + o(t - t_0)$，无法说明曲线
如何离开该点——这正是下文方法中对
奇点单独讨论的原因，其中第一个
\emph{非零}导数承担 $f'$ 的作用。
另外注意，切线是几何对象：任何
重新参数化仅使 $f'$ 乘以一个非零标量，并
不改变直线本身。
\end{proof}

\begin{method}[参数曲线的研究]\label{met:b1:curves:study}
\begin{enumerate}
  \item 利用对称性，\emph{缩小定义域}：考察
        $M(-t)$、$M(t + T)$、\dots\ 与 $M(t)$ 之间的关系（反射、
        平移、周期性），然后只画出缩小后的部分。
  \item \emph{变化情况}：联合列出 $x'$ 与 $y'$ 的符号表；
        曲线随 $x'$ 的符号向右或向左移动，随 $y'$ 的符号向上或向下移动。
  \item \emph{特殊点}：水平切线（$y' = 0 \neq
        x'$），竖直切线（$x' = 0 \neq y'$）；在
        \emph{奇点}（$f' = 0$）处，考察高阶导数
        以确定切线方向。
  \item \emph{渐近行为}：在 $I$ 的两端考察，然后作草图。
\end{enumerate}
\end{method}

\begin{example}[渐近分支，详解]
\label{ex:b1:curves:asymptotes}
$x(t) = t$，$y(t) = t + \dfrac1t$，定义于 $\intoo{0}{+\infty}$ 上。当 $t
\to 0^{+}$ 时：$x \to 0$，而 $y \to +\infty$——曲线沿
\emph{竖直渐近线} $x = 0$ 上升（一个\omterm{prop:b1:vspaces:coordinates}{坐标}有有限极限，
另一个趋于无穷）。当 $t \to +\infty$ 时：
两个\omterm{prop:b1:vspaces:coordinates}{坐标}都趋于无穷，故考察一条直线：差
\[
y(t) - x(t) = \frac1t \longrightarrow 0^{+}
\]
由此得出\emph{斜渐近线} $y = x$，且从
上方逼近。在两者之间，$y' = 1 - \frac1{t^2}$ 在 $t = 1$ 处为零：
点 $(1, 2)$ 为该支的最低点，且 $x' = 1 >
0$ 在全程成立，故曲线总是向右前进。一条弧，
两条渐近线，一个极小值：由三次计算
即得完整图形——其背后隐含一般方法：
若 $x, y \to \infty$ 同时成立，则考察 $y/x$ 以得候选
斜率（此处 $\to 1$），再考察 $y - (\text{斜率})\,x$ 以求
截距及逼近的一侧。
\end{example}

\begin{example}[星形线]\label{ex:b1:curves:astroid}
$x(t) = \cos^3 t$，$y(t) = \sin^3 t$。\index{星形线}对称性：
$M(t + 2\pi) = M(t)$（只研究一个周期）；$M(-t)$ 是
$M(t)$ 关于 $x$ 轴的对称点；$M(\pi - t)$ 关于 $y$ 轴的对称点；$M(\frac\pi2 -
t)$ 关于对角线 $y = x$ 的对称点：只需研究 $t \in
\intcc{0}{\frac\pi4}$ 上的部分，再展开即可。

速度：$f'(t) = 3\sin t\cos t\,(-\cos t, \sin t)$。在
$\intoo{0}{\frac\pi2}$ 上所有点均为正则点，切线方向由
$(-\cos t, \sin t)$ 确定；在 $t = 0$ 处（点 $(1,0)$），速度
为零：这是\emph{尖点}\index{尖点}，曲线在该处沿切线方向
$(-1, 0)\cdot(\pm)$ 折返——由对称性，四个
\omterm{ex:b1:curves:astroid}{尖点}位于 $(\pm1, 0), (0, \pm1)$。
\end{example}

\begin{example}[八字形曲线的完整研究]
\label{ex:b1:curves:lissajous}
$x(t) = \sin t$，$y(t) = \sin 2t$（\emph{利萨如曲线}）。
\emph{对称性}：$M(t + \pi) = (-x(t), y(t))$（关于
$y$ 轴的对称），$M(-t) = (-x(t), -y(t))$（中心对称），
$M(\pi - t) = (x(t), -y(t))$（关于 $x$ 轴的对称）：只需
研究 $t \in \intcc{0}{\frac\pi2}$ 的情形，再展开。
\emph{变化情况}：$x' = \cos t \geq 0$ 在全程成立，而 $y' =
2\cos 2t$ 在 $t = \frac\pi4$ 之前为正，之后为负：
曲线向右上方攀升至顶点
$\bigl(\frac{\sqrt2}2,\ 1\bigr)$，在 $t = \frac\pi4$ 处（水平
切线），然后向右下降至 $(1, 0)$，在 $t =
\frac\pi2$ 处，$x' = 0 \neq y'$：竖直切线。
\emph{二重点}：$M(0) = M(\pi) = (0,0)$，两次经过时
速度不同：
\[
f'(0) = (1,\ 2), \qquad f'(\pi) = (-1,\ 2) :
\]
两条\emph{正则}分支在原点以不同
角度相交——这是二重点，而非奇点：每次经过都是
完全光滑的，两次经过只是恰好共享同一位置。
整条曲线即下图所示的八字形。
\end{example}

\begin{omfigure}
\begin{tikzpicture}[scale=1.5]
  \draw[->, thin] (-1.25,0) -- (1.3,0) node[below] {$x$};
  \draw[->, thin] (0,-1.25) -- (0,1.3) node[left] {$y$};
  \draw[omDef, very thick, domain=0:360, samples=200, smooth]
    plot ({sin(\x)}, {sin(2*\x)});
  \fill (0.7071, 1) circle (0.035);
  \node[above, font=\small] at (0.7071, 1)
    {$\bigl(\tfrac{\sqrt2}2, 1\bigr)$};
  \fill (0,0) circle (0.035);
\end{tikzpicture}

{\small 利萨如曲线 $(\sin t, \sin 2t)$：原点处为一个二重点，
两条正则分支在此相交，速度分别为
$(1, 2)$ 与 $(-1, 2)$，并在四个
顶点处切线水平。利用对称性，全部计算只需在四分之一周期内进行。}
\end{omfigure}

\begin{omfigure}
\begin{tikzpicture}[scale=1.35]
  \draw[->, thin] (-1.3,0) -- (1.35,0) node[below] {$x$};
  \draw[->, thin] (0,-1.3) -- (0,1.35) node[left] {$y$};
  \draw[omDef, very thick, domain=0:360, samples=200, smooth]
    plot ({cos(\x)^3}, {sin(\x)^3});
  \node[below right, font=\small] at (1,0) {$1$};
  \node[above left, font=\small] at (0,1) {$1$};
\end{tikzpicture}

{\small 星形线 $(\cos^3 t, \sin^3 t)$：四段弧相交于
四个\omterm{ex:b1:curves:astroid}{尖点}。它是半径为
$\frac14$ 的圆在单位圆内滚动时，圆上一点所描出的轨迹。}
\end{omfigure}

\section{极坐标曲线}

\begin{definition}\label{def:b1:curves:polar}
称由 $r = r(\theta)$ 给出的曲线为\emph{极坐标曲线}\index{极坐标曲线}：
参数为 $\theta$ 的点为
\[
M(\theta) = r(\theta)\,\vec u(\theta),
\qquad \vec u(\theta) = (\cos\theta, \sin\theta).
\]
记 $\vec v(\theta) = (-\sin\theta, \cos\theta) = \vec
u\,'(\theta)$，则速度为
\[
M'(\theta) = r'(\theta)\, \vec u(\theta) + r(\theta)\, \vec
v(\theta) .
\]
由此可得：当 $r \neq 0$ 时，该点为正则点，
切线与射线所成的角 $V$ 满足 $\tan V =
\frac{r}{r'}$（即 $M'$ 与 $\vec u$ 的夹角）；当 $r(\theta_0)
= 0$ 时，曲线经过原点，\emph{切线为
射线 $\theta = \theta_0$}（方向为 $\vec u(\theta_0)$，由
$M'(\theta_0) = r'(\theta_0)\vec u(\theta_0)$ 读出，或取极限弦：
即从 $O$ 到 $M(\theta)$ 的弦沿 $\vec
u(\theta)$ 本身的方向，它趋于 $\vec u(\theta_0)$——因此
即使 $r'(\theta_0) = 0$ 且该点为奇点，该规则仍然成立，
这正是原点处的极\omterm{prop:b1:vspaces:coordinates}{坐标}\omterm{ex:b1:curves:astroid}{尖点}，如心形线那样，
无需额外工作即可得到切线。）
\end{definition}

\begin{example}[心形线]\label{ex:b1:curves:cardioid}
$r(\theta) = 1 + \cos\theta$。\index{心形线}对称性：$r(-\theta)
= r(\theta)$，关于 $x$ 轴对称；只需研究 $\theta \in
\intcc{0}{\pi}$。$r$ 从 $2$ 减小到 $0$；当 $\theta = \pi$ 时，
$r = 0$，曲线沿与射线相切的方向到达原点
$\theta = \pi$（即 $x$ 轴）相切，并在该处形成\omterm{ex:b1:curves:astroid}{尖点}——即心尖。
$\theta = 0$ 处：$r' = 0$，故 $\tan V = \infty$：
即切线与轴垂直。

关于角 $V$，还值得再读一遍。当 $\theta =
\frac\pi2$ 时，$r = 1$，$r' = -1$，故 $\tan V = \frac{r}{r'} =
-1$：切线与出射射线所成的角为
四分之三个直角——曲线正朝其\omterm{ex:b1:curves:astroid}{尖点}弯回。
公式 $\tan V = r/r'$ 给出整条曲线上的切线方向，
且完全无需计算 $M'(\theta)$，
这正是极\omterm{prop:b1:vspaces:coordinates}{坐标}下读取斜率的类比。
\end{example}

\begin{example}[极坐标化为直角坐标：隐藏的圆]
\label{ex:b1:curves:polarcircle}
\omterm{def:b1:curves:polar}{极坐标曲线} $r = 2\cos\theta$ 是什么？两边乘以 $r$：
$r^2 = 2r\cos\theta$，即 $x^2 + y^2 = 2x$，即
\[
(x - 1)^2 + y^2 = 1 :
\]
这是圆心为 $(1, 0)$、半径为 $1$ 的圆，且过
原点。这里的计数很重要：当 $\theta$ 遍历
$\intoc{-\frac\pi2}{\frac\pi2}$ 时，整个圆恰好被扫过
一次（$r$ 在两端都为零），且在 $\theta =
\pm\frac\pi2$ 时，规则“沿射线
$\theta = \theta_0$ 的原点切线”给出\emph{竖直}切线——
与几何事实相符，竖直轴确实与此圆在
$O$ 处相切。当 $\theta$ 超出该范围时，$r < 0$
重新描出同一圆，提醒我们：\omterm{def:b1:curves:polar}{极坐标曲线}是一个
\emph{参数化}对象，允许轨迹自我重叠。
\end{example}

\begin{omfigure}
\begin{tikzpicture}[scale=1.5]
  \draw[->, thin] (-1.2,0) -- (1.3,0) node[below] {$x$};
  \draw[->, thin] (0,-1.2) -- (0,1.3) node[left] {$y$};
  \draw[omProp, very thick, domain=0:360, samples=300, smooth]
    plot ({cos(2*\x)*cos(\x)}, {cos(2*\x)*sin(\x)});
  \draw[thin, dashed] (-0.85,-0.85) -- (0.85,0.85);
  \draw[thin, dashed] (-0.85,0.85) -- (0.85,-0.85);
\end{tikzpicture}

{\small 四瓣玫瑰线 $r = \cos 2\theta$
（见\cref{exo:b1:curves:4}）。沿 $y$ 轴方向的花瓣
由 $r < 0$ 描出（该点位于与 $\theta$ 相反的
射线上）；虚线对角线 $\theta = \pm\frac\pi4$ 即
原点处的切线，此处 $r$ 为零。}
\end{omfigure}

\begin{omfigure}
\begin{tikzpicture}[scale=1.15]
  \draw[->, thin] (-0.7,0) -- (2.4,0) node[below] {$x$};
  \draw[->, thin] (0,-1.7) -- (0,1.7) node[left] {$y$};
  \draw[omThm, very thick, domain=0:360, samples=200, smooth]
    plot ({(1+cos(\x))*cos(\x)}, {(1+cos(\x))*sin(\x)});
  \node[below, font=\small] at (2,0) {$2$};
\end{tikzpicture}

{\small 心形线 $r = 1 + \cos\theta$。\omterm{def:b1:curves:polar}{极坐标曲线}的读法是
让角度扫过：极径随 $\theta$
变化，先增大又减小。}
\end{omfigure}

\section{弧长}

\begin{definition}[弧长]\label{def:b1:curves:arclength}
\emph{弧长}\index{弧长}是 $C^1$ 弧 $f \colon
\intcc{a}{b} \to \R^2$ 的速率的\omterm{thm:b1:integration:def}{积分}：
\[
L = \int_a^b \norm{f'(t)}\,\dd t
= \int_a^b \sqrt{x'(t)^2 + y'(t)^2}\;\dd t .
\]
（动机：在小\omterm{prop:b1:reals:intervals}{区间}上 $M(t + h) \approx M(t) + h
f'(t)$，故弧近似于一条折线，其各段长之和
为 $\norm{f'}$ 的一个\omterm{thm:b1:integration:riemann}{黎曼和}，
见\cref{thm:b1:integration:riemann}。）对\omterm{def:b1:curves:polar}{极坐标曲线} $r =
r(\theta)$，速度 $r'\vec u + r\vec v$ 的两个分量正交，
因此
\[
L = \int_{\theta_1}^{\theta_2}
\sqrt{r'(\theta)^2 + r(\theta)^2}\;\dd\theta .
\]
弧长不依赖于所选的（单调、$C^1$）
参数化：在\omterm{thm:b1:integration:def}{积分}中作代换 $t = \varphi(s)$
（见\cref{thm:b1:integration:parts}），则 $f'$ 乘以
$\varphi'$，$\dd t$ 乘以 $\varphi'^{-1}$。
\end{definition}

\begin{example}[检验：圆]\label{ex:b1:curves:circlelength}
对 $f(t) = (R\cos t, R\sin t)$，在 $\intcc{0}{2\pi}$ 上：$\norm{f'} =
R$，故 $L = 2\pi R$——定义给出的正是圆周长。
重新参数化检验：\omterm{def:b1:logic:map}{映射} $g(t) = (R\cos 2t, R\sin 2t)$，在
$\intcc{0}{\pi}$ 上描出的是\emph{同一}个圆，速率加倍，
$\norm{g'} = 2R$，并且
\[
\int_0^{\pi} 2R\,\dd t = 2\pi R
\]
同样成立：时间减半、速率加倍、长度不变——这
正是定义中承诺的不变性，此处用数值验证了一次。
（若把 $g$ 用在整个 $\intcc{0}{2\pi}$ 上，得 $4\pi R$：
一条绕行两次的曲线，作为\emph{行程}确实长两倍；
长度度量的是带参数的路径，对\omterm{prop:b1:reals:intervals}{区间}的如实记录
本身就是计算的一部分。）
对星形线 $(\cos^3t, \sin^3t)$：$\norm{f'} = 3\abs{\sin
t\cos t} = \tfrac32\abs{\sin 2t}$，由对称性 $L = 4\int_0^{\pi/2}
\tfrac32\sin 2t\,\dd t = 6$：这是一条画在单位圆内的曲线，
其长度 $6 < 2\pi$。周末问题将测量
所有拱中最著名的那一拱。
\end{example}

\begin{omfigure}
\begin{tikzpicture}[scale=0.85]
  \draw[->, thin] (-0.5,0) -- (7.2,0) node[below] {$x$};
  \draw[->, thin] (0,-0.4) -- (0,2.6) node[left] {$y$};
  \draw[omThm, very thick, domain=0:6.2832, samples=160, smooth]
    plot ({\x - sin(\x r)}, {1 - cos(\x r)});
  \draw[omDef, thick] (2,1) circle (1);
  \fill[omDef] (2,1) circle (0.04);
  \fill (2 - 0.9093, 1.4161) circle (0.06);
  \draw[omProp, thick] (2,0) -- (2 - 0.9093, 1.4161) -- (2,2);
  \fill (2,0) circle (0.05);
  \fill (2,2) circle (0.05);
  \node[below, font=\small] at (2,0) {$C$};
  \node[above, font=\small] at (2,2) {$T$};
  \node[left, font=\small] at (2 - 0.9093, 1.4161) {$M$};
  \node[below, font=\small] at (6.2832,0) {$2\pi R$};
\end{tikzpicture}

{\small 摆线的一拱（半径 $R = 1$），
$t = 2$ 时的滚动圆，以及周末问题的两条辅助线：
到接触点的弦 $MC$ 是曲线的\emph{法向}，
到圆顶的弦 $MT$ 与曲线\emph{相切}。}
\end{omfigure}

\section{圆锥曲线}

\begin{definition}[焦点--准线定义]\label{def:b1:curves:conics}
固定一点 $F$（\emph{焦点}），一条直线 $D$
不经过 $F$（\emph{准线}）以及 $e > 0$（\emph{离心率}）。由这些数据确定的
\emph{圆锥曲线}\index{圆锥曲线}为
\[
\mathcal{C} = \{M : d(M, F) = e\; d(M, D)\}:
\]
\emph{椭圆}（$e < 1$）、\emph{抛物线}（$e = 1$）、
\emph{双曲线}（$e > 1$）。（圆可视为退化的
极限情形 $e \to 0$。）
\end{definition}

\begin{example}[在抛物线上检验定义]
\label{ex:b1:curves:parabolacheck}
取抛物线 $y^2 = 4x$，即 $2p = 4$：焦点 $F = (1, 0)$，
准线 $D : x = -1$（标准形式
由\cref{thm:b1:curves:reduced}，它们位于 $\pm\frac p2$ 处）。
在曲线上的点 $M = (1, 2)$：
\[
MF = \sqrt{(1-1)^2 + 2^2} = 2,
\qquad
d(M, D) = 1 - (-1) = 2 :
\]
两者相等，符合 $e = 1$ 的要求。在 $M' = (4, 4)$ 处：$MF' = \sqrt{9 +
16} = 5$，且 $d(M', D) = 5$，同样成立。焦点--准线
定义并非抽象之物：它是在任一点上都可以测量的一对距离，
而标准方程的代数
不过是将这一测量
一劳永逸地完成。
\end{example}

\begin{theorem}[标准方程]\label{thm:b1:curves:reduced}
在一组选取恰当的标准正交标架中：
\[
\text{椭圆： } \frac{x^2}{a^2} + \frac{y^2}{b^2} = 1
\quad (a \geq b > 0),
\qquad
\text{双曲线： } \frac{x^2}{a^2} - \frac{y^2}{b^2} = 1,
\qquad
\text{抛物线： } y^2 = 2px,
\]
其中对于椭圆：焦点位于 $(\pm c, 0)$，$c = \sqrt{a^2 - b^2}$，
$e = \frac ca$，以及双焦点刻画 $MF + MF' = 2a$；
对于双曲线：$c = \sqrt{a^2 + b^2}$，$e = \frac ca$，
$\abs{MF - MF'} = 2a$，渐近线 $y = \pm\frac ba x$。
\end{theorem}

\begin{proof}
取焦点为原点、准线竖直，即 $x = -h$
（$h > 0$），条件 $MF^2 = e^2 d(M, D)^2$ 化为 $x^2 + y^2 =
e^2 (x + h)^2$。当 $e = 1$ 时，$y^2 = 2hx + h^2$，
作平移 $x \mapsto x - \frac h2$ 后即为抛物线（$p = h$）。当 $e \neq 1$ 时，
对 $x$ 配方，
\[
(1 - e^2)\Bigl(x - \frac{e^2 h}{1 - e^2}\Bigr)^{\!2} + y^2
= e^2h^2 + \frac{e^4 h^2}{1 - e^2}
= \frac{e^2 h^2}{1 - e^2} .
\]
记 $X = x - \frac{e^2h}{1-e^2}$（\omterm{prop:b1:vspaces:coordinates}{坐标}平移）。当 $e < 1$ 时，
两边除以正的右端项：
得 $\frac{X^2}{a^2} + \frac{y^2}{b^2} = 1$，其中 $a = \frac{eh}{1 -
e^2}$，$b = \frac{eh}{\sqrt{1-e^2}}$；当 $e > 1$ 时，
等式两边适当变号，同样的除法给出
$\frac{X^2}{a^2} - \frac{y^2}{b^2} = 1$，其中 $a = \frac{eh}{e^2 -
1}$，$b = \frac{eh}{\sqrt{e^2 - 1}}$。所述 $c$ 的值
由 $c^2 = a^2 - b^2$ 或 $a^2 + b^2$ 均得 $c = ea$，
焦点位置由此正确；双焦点性质
可直接在化简后的方程上验证。对椭圆细述如下，
取 $F' = (c, 0)$，设 $M = (X, y)$ 为曲线上的点：
\[
MF'^2 = (X - c)^2 + y^2
= X^2 - 2cX + c^2 + b^2\Bigl(1 - \frac{X^2}{a^2}\Bigr)
= \frac{c^2}{a^2}X^2 - 2cX + a^2
= (a - eX)^2 ,
\]
利用 $c^2 + b^2 = a^2$ 与 $e = \frac ca$；由 $\abs X \leq
a$ 及 $e < 1$ 知 $a - eX > 0$，于是 $MF' = a - eX$，
全程无需开平方。镜像计算得 $MF = a + eX$，
从而 $MF + MF' = 2a$，为常数。对双曲线，
同样的代数运算给出 $MF = \abs{a + eX}$ 与 $MF' = \abs{eX - a}$，
差为 $\pm2a$，取决于分支。
\end{proof}

\begin{example}\label{ex:b1:curves:conicexample}
$\frac{x^2}{25} + \frac{y^2}{9} = 1$：椭圆，$a = 5$，$b = 3$，$c
= 4$：焦点 $(\pm4, 0)$，离心率 $\frac45$。其参数方程为：
$(5\cos t, 3\sin t)$——一个圆经不等比伸缩而成；
其上任一点到两焦点的距离之和为 $10$。
\end{example}

\begin{example}[由极坐标方程读出轨道]
\label{ex:b1:curves:orbit}
这条极\omterm{prop:b1:vspaces:coordinates}{坐标}\omterm{def:b1:curves:conics}{圆锥曲线} $r = \dfrac{1}{1 + \frac12\cos\theta}$
（见\cref{exo:b1:curves:7}，$p = 1$，$e = \frac12$）是一个
\emph{焦点在原点}的椭圆——行星轨道的几何，
即太阳位于 $O$ 的行星轨道。一切量都要从
$p$ 与 $e$ 出发，借助
\cref{thm:b1:curves:reduced}的化简，$a = \dfrac{p}{1 - e^2} =
\dfrac{1}{3/4} = \dfrac43$，$c = ea = \dfrac23$。两个拱点
可直接验证，不需任何理论：
\[
r(0) = \frac{1}{3/2} = \frac23 = a - c
\quad (\text{近日点}),
\qquad
r(\pi) = \frac{1}{1/2} = 2 = a + c
\quad (\text{远日点}),
\]
二者之和 $\frac23 + 2 = \frac83 = 2a$ 给出
长轴长。只要焦点在物理上居特殊地位，
极\omterm{prop:b1:vspaces:coordinates}{坐标}形式便最自然；只要对称轴更为重要，
化简后的直角\omterm{prop:b1:vspaces:coordinates}{坐标}形式便最自然。二者之间的互换，
恰是定理中配方计算所完成的事。

最后，把离心率看作一个旋钮：保持 $p = 1$，让 $e$ 在 $r
= \frac{1}{1 + e\cos\theta}$ 中变化：$e = 0$ 时为圆 $r = 1$；
当 $e = \frac12$ 时，为刚才研究的椭圆，$r$ 在
$\frac23$ 与 $2$ 之间振荡。当 $e = 1$ 时，$r(\theta) \to \infty$，当 $\theta
\to \pi$ 时——曲线不再闭合：为抛物线，其
最远点被推向无穷远。当 $e = 2$ 时，分母
在 $\cos\theta = -\frac12$ 处为零，仅 $\theta \in
\intoo{-\frac{2\pi}3}{\frac{2\pi}3}$ 保留，为双曲线的一支，
沿两个渐近方向远去。一个
公式涵盖整个\omterm{def:b1:curves:conics}{圆锥曲线}族，过渡点 $e =
1$ 则表现为分母首次为零的时刻。
\end{example}

\begin{remark}[常见陷阱]
\label{rem:b1:curves:pitfalls}
\emph{二重点不是奇点}：在自交点处，
每一支都是正则的；所谓“奇点”指的是
参数的某个取值使 $f'(t_0) = 0$，
（对照\cref{ex:b1:curves:lissajous}与星形线的\omterm{ex:b1:curves:astroid}{尖点}）。
\emph{竖直切线与\omterm{ex:b1:curves:astroid}{尖点}之别}：$x'(t_0) = 0 \neq y'(t_0)$
是具竖直切线的正则点；只有 $x' = y' = 0$
才需要更高阶的展开。
\emph{负半径画在反向射线上}：当 $r(\theta) <
0$ 时，该点为 $-\abs{r}\,\vec u(\theta)$，其极角为 $\theta +
\pi$（见\cref{exo:b1:curves:5}）。忽略此点则会漏掉内环
或使曲线重复出现。
\emph{\omterm{def:b1:curves:arclength}{弧长}\omterm{thm:b1:integration:def}{积分}的是 $\norm{f'}$，不是 $f'$}：应在速率的
零点处切开\omterm{thm:b1:integration:def}{积分}——对星形线，
在一个完整周期上不取绝对值地\omterm{thm:b1:integration:def}{积分} $\frac32\sin 2t$，
得 $0$，而不是 $6$。
\emph{\omterm{def:b1:curves:conics}{圆锥曲线}须先化简再读取}：在 $9x^2 +
25y^2 - 36x - 50y - 164 = 0$ 中，对称轴、
中心与离心率均看不出，直到配方完成
（见\cref{exo:b1:curves:6}）；若二次项在某一变量上消失，
则为抛物线，这不是“退化椭圆”。
\end{remark}

\begin{remark}[这些曲线的用途]
\label{rem:b1:curves:whereused}
\omterm{def:b1:curves:def}{参数曲线}是力学的语言：轨迹
是曲线，速度向量是切向量，
而\cref{def:b1:curves:arclength}所定义的\omterm{def:b1:curves:arclength}{弧长}即为所行的路程。
\omterm{def:b1:curves:conics}{圆锥曲线}在凡有平方反比律作用之处重新出现——行星
轨道即是太阳位于焦点的椭圆。在
\cref{ch:b1:multivar}中，曲线成为
二元函数的水平集，而本章的切线
与下一章的梯度相接。本科二年级分册增加了曲率与
曲线的局部标准形；下面的周末问题
仅用本年度的工具，即已推出十七世纪
人们对其最著名曲线所知的一切。
\end{remark}

\begin{remark}[第 3 册内的几种视角]
\label{rem:b1:curves:perspectives}
本章占去本册的整个前半部分，并
为末章提供素材。\emph{所用}：切向量为
导数（\cref{ch:b1:derivative}），\omterm{def:b1:curves:arclength}{弧长}与面积为
\omterm{thm:b1:integration:def}{积分}（\cref{ch:b1:integration}），\omterm{ex:b1:curves:astroid}{尖点}由
泰勒展开确定（\cref{ch:b1:taylor}），等时曲线为
线性微分方程（\cref{ch:b1:diffeq}），而每个
距离与角度皆属欧几里得几何（\cref{ch:b1:euclid}）。
\emph{所供}：在\cref{ch:b1:multivar}中，曲线 $t \mapsto (x(t),
y(t))$ 位于水平集 $\{f = c\}$ 之内，经链式法则求导，
所得恒等式 $\langle \nabla f,\
f'(t)\rangle = 0$ 将本章的切向量与下一章的梯度相接——
即曲线的切线与曲面的法线方向，
乃是同一计算
从两岸所见。
\end{remark}

\section{练习}

\begin{exercise}[$\star$]\label{exo:b1:curves:1}
研究并画出曲线 $x(t) = t^2$，$y(t) = t^3$（对称性、
变化情况、奇点 $t = 0$ 处的性态）。
\end{exercise}

\begin{exercise}[$\star$]\label{exo:b1:curves:2}
对于旋轮线 $x(t) = t - \sin t$，$y(t) = 1 - \cos t$（
即半径为 $1$ 的滚轮上一点的轨迹）：求
出周期平移对称性、奇点，以及
$t = 0$ 处的切线方向 \emph{（将 $x$ 与 $y$ 展开至
首个非零阶）}。
\end{exercise}

\begin{exercise}[$\star$]\label{exo:b1:curves:3}
求曲线 $(\cos^3 t, \sin^3 t)$ 在 $t =
\frac\pi4$ 处的切线，并验证该切线被\omterm{prop:b1:vspaces:coordinates}{坐标}轴所截线段
的长为 $1$ ——星形线的一个著名性质（在
每个正则点处皆成立）。
\end{exercise}

\begin{exercise}[$\star$]\label{exo:b1:curves:4}
画出\omterm{def:b1:curves:polar}{极坐标曲线} $r = \cos 2\theta$（四瓣玫瑰线）与
$r = \frac{1}{\cos\theta}$ 在 $\intoo{-\frac\pi2}{\frac\pi2}$ 上的图像
（并认出其为一条直线）。
\end{exercise}

\begin{exercise}[$\star\star$]\label{exo:b1:curves:5}
研究\omterm{def:b1:curves:polar}{极坐标曲线} $r = 1 + 2\cos\theta$（蜗线）：考察
$r \geq 0$ 与 $r < 0$ 所对应的范围（负极径的点位于
相反的射线上），过原点的情形、内圈，
并画出草图。
\end{exercise}

\begin{exercise}[$\star\star$]\label{exo:b1:curves:6}
识别\omterm{def:b1:curves:conics}{圆锥曲线} $9x^2 + 25y^2 - 36x - 50y - 164 = 0$：通过
配方化简，求出其中心、半轴、焦点与离心率。
\end{exercise}

\begin{exercise}[$\star\star$]\label{exo:b1:curves:7}
证明：以原点为焦点的\omterm{def:b1:curves:conics}{圆锥曲线}的极\omterm{prop:b1:vspaces:coordinates}{坐标}方程
为
\[
r = \frac{p}{1 + e\cos\theta}
\]
（准线为竖直直线，位于焦点右侧距离 $\frac pe$ 处）。
$\theta$ 在 $e > 1$ 时可取哪些值？
\end{exercise}

\begin{exercise}[$\star\star$]\label{exo:b1:curves:8}
由双焦点性质 $MF + MF' = 2a$，推出椭圆的“园丁作图法”，并证明
在 $M$ 处的切线
与 $MF$、$MF'$ 所成角相等 \emph{（反射性质；
对 $\norm{M(t) - F} + \norm{M(t) - F'} = 2a$ 求导，
并解释单位向量乘积之和为零）}。
\end{exercise}

\begin{exercise}[$\star\star\star$]\label{exo:b1:curves:9}
\emph{伯努利双纽线}是\omterm{def:b1:curves:polar}{极坐标曲线} $r^2 = \cos
2\theta$（在其有定义处取 $r = \sqrt{\cos2\theta}$）。
\begin{enumerate}
  \item 求其定义域、对称性、原点处的切线，并
        画出图形。
  \item 证明其为点 $M$ 的轨迹，这些点满足 $MF \cdot MF'
        = \frac12$，其中 $F, F' = \bigl(\pm\frac{1}{\sqrt2},
        0\bigr)$。\emph{（用极\omterm{prop:b1:vspaces:coordinates}{坐标}计算
        $MF^2\,MF'^2$。）}
\end{enumerate}
\end{exercise}

\begin{exercise}[$\star\star$]\label{exo:b1:curves:10}
计算心形线 $r = 1 + \cos\theta$ 的总长度，
用\cref{def:b1:curves:arclength}中的极\omterm{prop:b1:vspaces:coordinates}{坐标}\omterm{def:b1:curves:arclength}{弧长}公式，
\emph{（恒等式 $1 + \cos\theta = 2\cos^2\frac\theta2$ 将平方根
化为 $2\abs{\cos\frac\theta2}$）}。
\end{exercise}

\begin{exercise}[$\star\star$]\label{exo:b1:curves:11}
证明椭圆 $(a\cos t, b\sin t)$ 在
参数为 $t$ 的点处的切线方程为
\[
\frac{x\cos t}{a} + \frac{y\sin t}{b} = 1 ,
\]
并由此推出椭圆 $\frac{x^2}{a^2} + \frac{y^2}{b^2} = 1$
在其上一点 $(x_0, y_0)$ 处的切线为 $\frac{x\,x_0}{a^2} +
\frac{y\,y_0}{b^2} = 1$（即“平方拆分”法则）。
\end{exercise}

\begin{exercise}[$\star\star\star$]\label{exo:b1:curves:12}
\emph{对数螺线}是\omterm{def:b1:curves:polar}{极坐标曲线} $r =
\eu^{k\theta}$（$k > 0$ 固定，$\theta \in \R$）。
\begin{enumerate}
  \item 证明极径与切线的夹角 $V$
        为\emph{常数}（$\tan V = \frac1k$）——
        即螺线与每条射线所成角都相同。
  \item 证明将螺线旋转角 $c$ 后，恰与其
        在因子 $\eu^{-kc}$ 下的伸缩像重合：每一次
        旋转都是螺线本身的一次伸缩
        （自相似性）。
  \item 计算弧段 $\theta \in
        \intoc{-\infty}{\theta_0}$ 的长度（作为
        $\intcc{A}{\theta_0}$ 上长度的极限，$A \to -\infty$），并验证其
        长度有限：这是一条绕原点盘旋无穷多圈的曲线，
        而其长度仍为有限。
\end{enumerate}
\end{exercise}

\section{问题：摆线，曲线中的女王}

\begin{problem}\label{pb:b1:curves:1}
半径为 $R$ 的车轮沿 $x$ 轴无滑动地滚动；
轮缘上初始位于原点的点所描出的曲线是
\emph{摆线}。十七世纪，学者们为这条曲线争论不休
——伽利略曾称量它的纸片剪影，雷恩测量过它的长度，
罗贝瓦尔算出了它的面积，惠更斯以它为原理造出时钟——而
他们得到的每一个结果都在本章的范围之内。
本章证明四个经典结论：切线作图、雷恩的长度
$8R$、面积 $3\pi R^2$ 以及惠更斯的等时
性质。\index{摆线}\index{等时曲线}

\medskip
\noindent\textbf{第一部分——滚动与切线。}
\begin{enumerate}
  \item 车轮转过角 $t$ 后，其中心
        位于 $\Omega(t) = (Rt, R)$（无滑动滚动：
        接触点的距离 $=$ 滚过的\omterm{def:b1:curves:arclength}{弧长}）。证明标记点
        位于
        \[
        M(t) = \bigl(R(t - \sin t),\; R(1 - \cos t)\bigr).
        \]
  \item 计算 $f'(t)$，并证明 $\norm{f'(t)} =
        2R\,\abs{\sin\frac t2}$；求出奇点
        （\omterm{ex:b1:curves:astroid}{尖点}——参见\cref{exo:b1:curves:2}）以及
        每一拱的拱顶。
  \item 设 $C(t) = (Rt, 0)$ 为接触点，$T(t) =
        (Rt, 2R)$ 为车轮的最高点。证明当 $0 < t <
        2\pi$ 时，向量 $M - C$ 为曲线在
        $M(t)$ 处的\emph{法向量}，而向量 $T - M$ 为\emph{切向量}：要作
        摆线的切线，只需将该点与其滚动圆的
        最高点相连。\emph{（一切都用
        $\sin\frac t2$ 与 $\cos\frac t2$ 表示。）}
  \item 从运动学角度解释第 3 小题：接触点是
        \emph{瞬时转动中心}，而
        $M$ 的速率等于它到 $C$ 的距离（取单位角速度），
        并验证 $\norm{M - C} = 2R\abs{\sin\frac t2}$。
  \item 证明高度--速率关系
        \[
        \norm{f'(t)}^2 = 2R\,y(t) :
        \]
        对以单位角速度运动的摆线，
        每一点的速率恰好等于
        自由下落高度为当前高度时的速率。（留待
        第四部分使用。）
\end{enumerate}

\medskip
\noindent\textbf{第二部分——雷恩定理：拱的长度为
$8R$。}
\begin{enumerate}[resume]
  \item 利用\cref{def:b1:curves:arclength}，计算
        一拱的长度：
        \[
        L = \int_0^{2\pi} 2R\sin\frac t2\,\dd t = 8R
        \]
        （雷恩定理，1658 年）。恰为四个车轮直径，
        且不出现任何 $\pi$。
  \item 计算从\omterm{ex:b1:curves:astroid}{尖点}起算的\omterm{def:b1:curves:arclength}{弧长}：$s(t) =
        4R\bigl(1 - \cos\frac t2\bigr)$，并验证 $s(2\pi) =
        8R$。
  \item 现在从\emph{拱顶}（$t = \pi$）起测量\omterm{def:b1:curves:arclength}{弧长}：
        $\sigma(t) = \abs{4R\cos\frac t2}$。证明
        内蕴关系
        \[
        \sigma^2 = 8R\,\bigl(2R - y\bigr) :
        \]
        从拱顶起算的\omterm{def:b1:curves:arclength}{弧长}的平方，与
        低于拱顶的高度降成正比。
  \item 检验：由第 8 小题可知，从拱顶到\omterm{ex:b1:curves:astroid}{尖点}的半拱
        长度为 $4R$，并与
        \cref{ex:b1:curves:circlelength}中的星形线计算相比较
——两条曲线的长度都是不含 $\pi$ 的\omterm{pb:b1:logic:1}{代数数}；
        试说明是什么使这一点成为可能，尽管二者
        都由圆生成。\emph{（观察
        $\norm{f'}$ 的形式。）}
\end{enumerate}

\medskip
\noindent\textbf{第三部分——罗贝瓦尔的面积：$3\pi R^2$。}
\begin{enumerate}[resume]
  \item 说明一拱与地面所围的面积为
        $A = \int_0^{2\pi} y(t)\,x'(t)\,\dd t$ \emph{（即
        将 $x = x(t)$ 代入 $\int y\,\dd x$，
        见\cref{thm:b1:integration:parts}；$x$ 递增）}。
  \item 计算
        \[
        A = R^2\int_0^{2\pi}(1 - \cos t)^2\,\dd t = 3\pi R^2 :
        \]
        恰为车轮面积的\emph{三}倍——
        即伽利略曾借助称重猜出的那个比值。
  \item 用同样的方法处理星形线 $(\cos^3t, \sin^3t)$，证明
        所围面积为 $\frac{3\pi}8$
        \emph{（将 $\sin^4 t\cos^2 t$ 线性化；在完整周期上，
        只有常数项保留）}。
  \item 在上半单位圆上检验第 10 小题的公式：
        $(\cos t, \sin t)$，$t$ 从 $\pi$ 变到 $0$：
        它是否给出 $\frac\pi2$？
\end{enumerate}

\medskip
\noindent\textbf{第四部分——惠更斯的等时曲线。}将拱线
        翻转过来：一颗无摩擦的小珠在重力 $g$ 作用下，沿
        摆线形碗的内侧滑动，其参数方程为
\[
x(t) = R(t + \sin t), \qquad y(t) = R(1 - \cos t)
\qquad (t \in \intcc{-\pi}{\pi}),
\]
        其最低点为原点（$y$ 向上计量）。
\begin{enumerate}[resume]
  \item 计算速率 $\norm{f'(t)} = 2R\cos\frac t2$，
        自最低点起算的\omterm{def:b1:curves:arclength}{弧长} $s(t) = 4R\sin\frac t2$，并
        证明关键恒等式
        \[
        y = \frac{s^2}{8R} .
        \]
  \item 珠子自参数为
        $t_0 > 0$ 的点静止释放，能量守恒：若 $s(\tau)$ 表示
        时刻 $\tau$ 的\omterm{def:b1:curves:arclength}{弧长}位置，则 $\frac12\bigl(
        \frac{\dd s}{\dd\tau}\bigr)^2 + g\,y = g\,y_0$。改写
        此式，利用第 14 小题，写为
        \[
        \Bigl(\frac{\dd s}{\dd\tau}\Bigr)^{\!2}
        = \frac{g}{4R}\,\bigl(s_0^2 - s^2\bigr),
        \qquad s_0 = s(t_0).
        \]
  \item 对 $\tau$ 求导，得到
        调和振子方程
        \[
        \frac{\dd^2 s}{\dd\tau^2} = -\frac{g}{4R}\,s ;
        \]
        利用\cref{thm:b1:diffeq:homogeneous2}求解，
        并由初值条件得 $s(\tau) = s_0\cos(\omega\tau)$，其中
        $\omega = \sqrt{g/4R}$。
  \item 由此得到\emph{等时性质}（惠更斯，1659 年）：
        到达最低点的时间，
        \[
        T_{\downarrow} = \frac{\pi}{2\omega}
        = \pi\sqrt{\frac Rg}\,,
        \]
        与释放点无关——从碗中任意两个高度
        同时释放的小珠，将
        同时到达最低点。
  \item 代入验证 $s(\tau) =
        s_0\cos\omega\tau$ 精确满足第 15 小题的一阶能量方程
        （而不只是求导后的那一式），
        并用一句话说明为何圆摆只是
        \emph{近似}等时的，
        而摆线则严格等时。
  \item 数值计算：何种半径 $R$ 使下降时间
        恰为一秒（$g = 9.81$）？注意答案距一米
        有多近，并比较完整振动
        周期 $2\pi\sqrt{4R/g}$ 与小角度摆公式
        $2\pi\sqrt{\ell/g}$（取 $\ell = 4R$）。
\end{enumerate}

\medskip
\noindent\textbf{第五部分——收获与综合。}
\begin{enumerate}[resume]
  \item 对第 3 小题的切线法则，在 $t =
        \frac\pi2$ 处应用（取 $R = 1$），计算 $M$、$T$、
        $MT$ 的方向，并将其
        与 $f'(\frac\pi2)$ 作比较。
  \item （旋轮线）改取到圆心距离为 $d$ 的点（$d \neq R$），
        此时曲线为 $x = Rt -
        d\sin t$，$y = R - d\cos t$。证明：当 $d < R$ 时，
        曲线处处正则，且不存在奇异行为（$x' > 0$：
        曲线总向前推进），而当 $d
        > R$ 时，横\omterm{prop:b1:vspaces:coordinates}{坐标} $x'$ 变号，曲线出现
        圈环。这正是轮缘上的点
        \emph{向后}运动的带凸缘铁路车轮。
  \item （最速降线预告）从碗的\omterm{ex:b1:curves:astroid}{尖点} $(\pi R, 2R)$
        出发至碗底，比较摆线的下降时间
        $\pi\sqrt{R/g}$ 与沿连接同两点的直线滑道下滑的时间
        \emph{（沿弦的加速度恒为
        $g\sin\alpha$）}：证明沿弦下滑所需时间为
        $\sqrt{\pi^2 + 4}\,\sqrt{R/g} \approx 3.72\sqrt{R/g}$。
        曲线胜过直线——事实上，它是所有曲线中最快的，
        这是变分法的结果。
  \item 证明在拱（$0 < t < 2\pi$）上，
        \[
        \frac{\dd y}{\dd x} = \cot\frac t2,
        \qquad
        \frac{\dd^2 y}{\dd x^2} =
        -\frac{1}{4R\sin^4\frac t2} < 0 :
        \]
        拱是凹的，且切线恰在
        \omterm{ex:b1:curves:astroid}{尖点}处竖直。
  \item 用几何方法重新得到
        \cref{exo:b1:curves:2} 中的竖直\omterm{ex:b1:curves:astroid}{尖点}切线：求
        第 3 小题中弦 $MT$ 当 $t \to
        0^{+}$ 时的极限方向，不作任何展开。
  \item 综合题，用四句话回答：本题证明了哪三个具名定理
        （附其结论 $8R$、$3\pi
        R^2$、$\pi\sqrt{R/g}$ 及其作者）；哪一种计算手法
        （$\sin\frac t2$、$\cos\frac t2$
        的因式分解）贯穿了第 I、II 与 IV 部分；
        内蕴关系 $y = s^2/8R$ 如何把几何
        化为线性微分方程；以及
        各部分分别依赖本书的哪一章。
\end{enumerate}
\end{problem}
